\chapter{Capa de aplicaciones}
\label{sec:capa-apps}

Las aplicaciones son la piel visible del sistema. Son tres:
\filepath{apps/tauri} (GUI de escritorio), \filepath{apps/cli}
(línea de comandos) y el servidor MCP TypeScript
\filepath{platform/mcp}. Las tres dependen de
\filepath{platform/api} como cliente; ninguna reimplementa lógica
propia.

\section{\filepath{apps/tauri}}

El desktop app. Su propósito es presentar el grafo, correr hunts
interactivos, mostrar scoring, y permitir anotación. Estructura
interna:

\begin{itemize}
    \item \filepath{src-tauri/src/commands/}: handlers de IPC
    Tauri, uno por dominio funcional (ingest, hunt, export,
    session, \ldots). Cada comando es un \textbf{pass-through}: recibe
    el payload, llama al método correspondiente de
    \filepath{platform/api}, devuelve el resultado.
    \item \filepath{src-tauri/src/http_api.rs}: routers Axum que
    exponen los 55 endpoints HTTP. Sirven tanto al MCP server como
    a scripts externos. Estos routers no tienen lógica; son
    adaptadores de transporte.
    \item \filepath{src/}: frontend React. Cytoscape para el grafo,
    componentes por vista (Hunt, Explore, Datasets, Ingest, Notes,
    Export). Context API, sin estado global pesado.
\end{itemize}

La consigna aquí es \textbf{YAGNI agresivo}: ningún componente del
Tauri debe anticipar variación que todavía no existe. Si dos
comandos comparten boilerplate, se factoriza a un helper, no a un
trait con una sola implementación.

\section{\filepath{apps/cli}}

El CLI es un consumidor in-proc de \filepath{platform/api}. No
corre un servidor HTTP; crea una instancia de
\code{GraphHunterApi}, llama operaciones y escribe a stdout en
formato legible o JSON.

Tres usos principales:

\begin{itemize}
    \item Ingesta batch de archivos (EVTX, CSV) a un workspace
    nuevo.
    \item Ejecutar un hunt desde un archivo de hipótesis y volcar
    resultados.
    \item Interactuar con \filepath{platform/siem} para trabajos
    programados de streaming (cron o scheduler externo).
\end{itemize}

El CLI tiene \textbf{cero} dependencia nueva sobre
\filepath{core/} desde la Fase~2: todo pasa por la API. Esto dejó
obsoleta a la antigua ``ruta paralela'' del CLI que bifurcaba la
lógica de ingesta.

\section{\filepath{platform/mcp}}

El servidor MCP (Model Context Protocol) es un paquete TypeScript
que expone 54 tools organizadas en 10 dominios, consumibles por
clientes de IA como Claude Code o Cursor.

Arquitectura (post ADR-004):

\begin{itemize}
    \item \filepath{src/lib/types.ts}: interfaz \code{Tool<I, O>}
    tipada, con \code{stability} explícita y \code{version}.
    \item \filepath{src/lib/registry.ts}: registro dinámico,
    agrupado por dominio.
    \item \filepath{src/tools/<dominio>/<tool>.ts}: una tool por
    archivo. Cada archivo expone un objeto que cumple el contrato
    \code{Tool<I, O>}.
\end{itemize}

Los 10 dominios son: \code{ingest}, \code{hunt}, \code{explore},
\code{export}, \code{session}, \code{agentic}, \code{dlq},
\code{invariants}, \code{mapping\_library}, \code{sentinel}.

Cada tool es un cliente HTTP contra \filepath{platform/api}. No
hay lógica de negocio en el MCP; cambios en el ranking de hunts,
por ejemplo, viven en \filepath{platform/api} una sola vez y se
propagan automáticamente.

\section{Por qué el MCP es plataforma, no aplicación}

Aunque el paquete vive físicamente en \filepath{platform/mcp/} y
nosotros lo tratamos como aplicación en el diagrama, la
clasificación oficial es \emph{plataforma}. La razón: el MCP
define \textbf{un contrato público} visto por clientes externos
(Claude, Cursor, custom agents). Cambios en su superficie afectan
consumidores que no controlamos, igual que cambios en los 55
endpoints HTTP. Por eso lleva el rigor de plataforma: versionado
por tool, tests contractuales, gate de estabilidad.

\section{Convenciones de aplicaciones}

Tres reglas operacionales para este nivel:

\begin{enumerate}
    \item \textbf{Delgado}: si un comando Tauri o una tool MCP
    contiene más de, digamos, quince líneas de lógica real, hay que
    mover esa lógica a \filepath{platform/api} y volver a dejar la
    superficie delgada.
    \item \textbf{Sin tests unitarios pesados}: smoke end-to-end
    sí; tests con mocks para cada handler, no. Si la lógica merece
    test unitario, pertenece a \filepath{platform/api} o al núcleo.
    \item \textbf{Sin traits preventivos}: no se introducen traits
    en \filepath{apps/} salvo que haya dos implementaciones reales
    \emph{hoy}. Esto contrasta con el estilo de
    \filepath{platform/}, y es deliberado.
\end{enumerate}
